%% kapa: exam=2 round=25 year=2014 period=3 subject=compensation_law subject_ko=감정평가 및 보상법규 seq=1 number=1 points=40
%% source: Q-Net 감정평가사 기출문제 (cst00309) artlSeq=5109209 file=extracted/1008/5109209/2014년 제25회 감정평가사 2차시험 기출문제/3교시(감정평가 및 보상법규).hwp
%% conversion: kordoc 4.24.51 -> tools/gen2.py (auto); review=proofread-sonnet; flags=none
\begin{longquestion}{【문제 1】}
S시의 시장 A는 K구의 D지역(주거지역)을「도시 및 주거환경정비법」(이하 “도정법”이라 함)상 정비구역으로 지정ㆍ고시하였다. 그러자 이 지역의 주민들은 조합을 설립하여 주택재개발사업을 추진하기 위해 도정법에서 정한 절차에 따라 조합설립추진위원회를 구성하였고, 동 추진위원회는 도정법 제16조의 규정에 의거하여 D지역의 일정한 토지등소유자의 동의, 정관, 공사비 등 정비사업에 드는 비용과 관련된 자료 등을 첨부하여 A로부터 X조합설립인가를 받아 등기하였다. X조합은 조합총회를 개최하고 법 소정의 소유자 동의 등을 얻어 지정개발자로서 Y를 사업시행자로 지정하였다. 다음 물음에 답하시오. (40점) \par
1) D지역의 토지소유자 중 甲이 “추진위원회가 주민의 동의를 얻어 X조합을 설립하는 과정에서 ‘건설되는 건축물의 설계의 개요’ 등에 관한 항목 내용의 기재가 누락되었음에도 이를 유효한 동의로 처리하여 조합설립행위에 하자가 있다.”고 주장하며 행정소송으로 다투려고 한다. 이 경우 조합설립인가의 법적 성질을 검토한 다음, 이에 기초하여 쟁송의 형태에 대해 설명하시오. (20점) \par
2) Y는 정비사업을 실시함에 있어 이 사업에 반대하는 토지등소유자 乙등의 토지와 주택을 취득하기 위하여「공익사업을 위한 토지 등의 취득 및 보상에 관한 법률」에 의거한 乙등과 협의가 성립되지 않아 지방토지수용위원회의 수용재결을 거쳤는데, 이 수용재결에 불복하여 Y가 중앙토지수용위원회에 이의재결을 신청하여 인용재결을 받았다. 이 경우 乙등이 이 재결에 대해 항고소송을 제기한다면 소송의 대상은 무엇인가? (20점) \par
\end{longquestion}
